% ============================================================
% Come si usano, argomento 26: codifica di sorgente
% ============================================================
\argomento{La codifica di sorgente}{Formule a p.~\pageref*{f:codice}}

\begin{scelta}{Come scegliere la formula}
\renewcommand{\arraystretch}{1.25}
\begin{tabular}{@{}P{0.55\linewidth}P{0.39\linewidth}@{}}
  \toprule
  \textbf{Se il testo\ldots} & \textbf{si usa\ldots} \\
  \midrule
  chiede la lunghezza media di un codice & \vv{codice} \\
  chiede se un codice si decifra in un solo modo & \vv{prefisso}, \vv{kraft} \\
  chiede se esiste un codice a prefisso con certe lunghezze & la disuguaglianza di Kraft \vv{kraft} \\
  chiede quanto si può comprimere la sorgente & il limite entropico \vv{limiteentropico} \\
  chiede il codice di Shannon, o il codice ottimo & \vv{shannoncodice}, \vv{huffman} \\
  chiede come avvicinarsi all'entropia quando le probabilità non sono diadiche & \vv{blocchi} \\
  \bottomrule
\end{tabular}

Prima si calcola l'entropia (\vv{entropia}): è il traguardo. Ogni codice decifrabile ha $L \geq H$, il codice di
Huffman è il migliore, quello di Shannon resta sotto $H + 1$.
\end{scelta}

\scheda{codice}{
  \usoquando{Si chiede quanti bit per simbolo usa, in media, un codice.}
  \usocome{Si moltiplica la lunghezza di ogni parola per la probabilità della sua lettera, e si somma.}
  \usoesempio{Sorgente $\left( \frac{1}{2}, \frac{1}{4}, \frac{1}{8}, \frac{1}{8} \right)$ e codice \texttt{0},
    \texttt{10}, \texttt{110}, \texttt{111}: $L = 1 \cdot \frac{1}{2} + 2 \cdot \frac{1}{4} + 3 \cdot \frac{1}{8} + 3
    \cdot \frac{1}{8} = 1{,}75$ bit/simbolo, contro i 2 del codice a lunghezza fissa.}
}

\scheda{prefisso}{
  \usoquando{Si chiede se un codice a lunghezza variabile si può decifrare senza ambiguità.}
  \usocome{Si controlla, per ogni coppia di parole, che la più corta non sia l'inizio della più lunga; se il controllo
    fallisce, si cerca una stringa che si legga in due modi.}
  \usoesempio{Il codice \texttt{0}, \texttt{1001}, \texttt{100}, \texttt{1111}, \texttt{01}, \texttt{110},
    \texttt{1100}, \texttt{111} non è decifrabile: \texttt{1111000} si legge \texttt{111|100|0}, cioè $(8, 3, 1)$, ma
    anche \texttt{1111|0|0|0}, cioè $(4, 1, 1, 1)$.}
}

\scheda{kraft}{
  \usoquando{Si chiede se esiste un codice a prefisso con certe lunghezze, o se un codice può essere decifrabile.}
  \usocome{Si somma $2^{-l}$ su tutte le parole: se la somma supera 1, nessun codice decifrabile ha quelle lunghezze.}
  \usoesempio{Lunghezze $2, 3, 3, 4, 3, 3, 4, 3$: $\frac{1}{4} + \frac{5}{8} + \frac{2}{16} = 1$, albero pieno.
    Lunghezze $1, 4, 3, 4, 2, 3, 4, 3$: la somma è $1{,}3125 > 1$, impossibile.}
}

\scheda{limiteentropico}{
  \usoquando{Si chiede di quanto si può comprimere, al massimo, una sorgente, o si deve giudicare un codice.}
  \usocome{Si calcola l'entropia e la si confronta con la lunghezza media: un codice con $L < H$ non può essere
    decifrabile.}
  \usoesempio{Sorgente a otto lettere con $H \simeq 2{,}80$ bit: il codice con $L = 2{,}5$ dell'esempio precedente
    scende sotto l'entropia, e infatti non è decifrabile; quello a prefisso ha $L \simeq 2{,}84$.}
}

\scheda{shannoncodice}{
  \usoquando{Si vuole un codice a prefisso costruito direttamente dalle probabilità, senza algoritmi.}
  \usocome{Per ogni lettera si calcola l'informazione in bit e la si arrotonda per eccesso (\vv{parteintera},
    \vv{logbase2}).}
  \usoesempio{Probabilità $\frac{1}{3}, \frac{1}{9}, \frac{1}{9}, \frac{1}{9}, \frac{1}{27}, \frac{2}{27},
    \frac{1}{6}, \frac{1}{18}$: lunghezze $2, 4, 4, 4, 5, 4, 3, 5$ e $L \simeq 3{,}26$, contro $H \simeq 2{,}70$.}
}

\scheda{huffman}{
  \usoquando{Si chiede il codice a prefisso più corto possibile.}
  \usocome{Conviene scrivere le probabilità con un denominatore comune; a ogni passo si uniscono le due più piccole e si
    tiene conto di quante volte ciascuna lettera viene unita: è la lunghezza della sua parola.}
  \usoesempio{Le probabilità di prima, in cinquantaquattresimi: $18, 6, 6, 6, 2, 4, 9, 3$. Lunghezze $2, 3, 3, 3, 4, 3,
    3, 4$ e $L = \frac{149}{54} \simeq 2{,}76$ bit, mezzo bit meno del codice di Shannon.}
}

\scheda{blocchi}{
  \usoquando{Le probabilità non sono diadiche e si vuole scendere quasi fino all'entropia.}
  \usocome{Si trattano i blocchi di $k$ lettere come lettere di un alfabeto più grande e li si codifica: lo spreco di
    al più un bit si divide fra le $k$ lettere.}
  \usoesempio{Moneta con $p = 0{,}11$, $H_2 \simeq 0{,}5$ bit: lettera per lettera serve 1 bit; con blocchi di 10
    lanci bastano meno di $0{,}6$ bit per lancio.}
}
